\section*{NeurIPS Paper Checklist}

\begin{enumerate}

\item {\bf Claims}
    \item[] Question: Do the main claims made in the abstract and introduction accurately reflect the paper's contributions and scope?
    \item[] Answer: \answerYes{}
    \item[] Justification: The abstract and Section~\ref{sec:intro}
    state exactly the four claims validated by the experiments
    (recovery, unbiasedness, concentration, and wall-clock speedup),
    with the downstream DiCo integration flagged separately, and each
    claim is numbered and cross-referenced. The scope restrictions
    (IPPO only, VMAS Multi-Agent Goal Navigation and Dispersion, fixed
    Gaussian policies; the main DiCo experiment is at a single set
    point $\SND_{\mathrm{des}}{=}0.1$, with a three-seed, three-set-point
    sweep at $n{=}50$ in Appendix~\ref{app:dico-n50-sweep}) are stated in
    Section~\ref{sec:discussion} (Limitations) and do not contradict
    any claim in the abstract or introduction.

\item {\bf Limitations}
    \item[] Question: Does the paper discuss the limitations of the work performed by the authors?
    \item[] Answer: \answerYes{}
    \item[] Justification: Section~\ref{sec:discussion} contains a
    dedicated ``Limitations'' paragraph itemising five scope limits
    (scale/timing regime coverage, scenario and DiCo-configuration
    coverage, fixed-$G$ ablation, closed-form Wasserstein assumption,
    IPPO-only scope).

\item {\bf Theory assumptions and proofs}
    \item[] Question: For each theoretical result, does the paper provide the full set of assumptions and a complete (and correct) proof?
    \item[] Answer: \answerYes{}
    \item[] Justification: All seven formal results
    (Propositions~\ref{prop:recovery}--\ref{prop:unbiased},
    Theorem~\ref{thm:concentration}, and Theorem~\ref{thm:distortion})
    are proved in full in Appendix~\ref{app:proofs}.
    Propositions~\ref{prop:recovery}--\ref{prop:unbiased} and
    Theorem~\ref{thm:concentration} are stated with their assumptions in
    Sections~\ref{sec:method}--\ref{sec:theory}; Theorem~\ref{thm:distortion}
    is stated in Appendix~\ref{app:proofs} with a compact summary in
    Section~\ref{sec:theory}. Assumptions (e.g.,
    non-negative edge weights with $W(G)>0$, bounded distances
    $d(i,j)\in[0,D_{\max}]$, independent Bernoulli inclusion under
    Horvitz--Thompson weights) are stated inside the corresponding
    theorem statements and reused verbatim in the proofs.

\item {\bf Experimental result reproducibility}
    \item[] Question: Does the paper fully disclose all the information needed to reproduce the main experimental results of the paper?
    \item[] Answer: \answerYes{}
    \item[] Justification: Section~\ref{sec:experiments:setup}
    specifies the VMAS environment, team sizes, policy architecture,
    training algorithm (IPPO), checkpoints (iter~0, iter~100), rollout
    set size, and metric implementation. Appendix~\ref{app:hyperparams}
    lists every PPO hyperparameter and the $n{=}100$ scaling overrides.
    The frozen-init GPU timing sweep and the DiCo integration
    (Sections~\ref{sec:experiments:timing},
    \ref{sec:experiments:dico-dispersion}) specify batch
    shape, timing regime (\texttt{torch.cuda.synchronize} on both
    sides of every \texttt{perf\_counter}), number of trials, number
    of seeds, and configuration of the $k$-NN graph and the Bernoulli
    random graph. The $n{=}50$ set-point sweep
    (Appendix~\ref{app:dico-n50-sweep}) reuses these hyperparameters
    except for the reduced per-iteration rollout
    (\texttt{on\_policy\_n\_envs\_per\_worker}${=}120$,
    \texttt{on\_policy\_collected\_frames\_per\_batch}${=}12{,}000$)
    and is driven by
    \texttt{ControllingBehavioralDiversity-fork/scripts/launch\_n50\_bern\_setpoint\_sweep.sh}.
    Everything needed to reconstruct the code paths is
    in \texttt{graphsnd/metrics.py}, \texttt{graphs.py},
    \texttt{training/train\_navigation\_batched.py}, and
    \texttt{experiments/exp2\_timing\_scaling.py} in the companion
    repository.

\item {\bf Open access to data and code}
    \item[] Question: Does the paper provide open access to the data and code, with sufficient instructions to faithfully reproduce the main experimental results, as described in supplemental material?
    \item[] Answer: \answerYes{}
    \item[] Justification: All code is released at the repository
    cited in the Conclusion, with a top-level \texttt{README.md}
    that specifies installation, training, experiment reproduction
    commands, and seed-42 expected values for each of the four
    validated claims. Raw CSVs, summary JSONs, and figure PDFs for
    every experiment in the paper are checked into the repository
    under \texttt{results/}. A preflight script
    (\texttt{scripts/preflight.sh}) runs the test suite and a
    smoke-training pass on each visible CUDA device to catch
    environment problems before overnight runs. For the anonymous
    submission the repository URL is anonymised in the main text;
    the supplementary material contains an anonymised snapshot.

\item {\bf Experimental setting/details}
    \item[] Question: Does the paper specify all the training and test details necessary to understand the results?
    \item[] Answer: \answerYes{}
    \item[] Justification: Training details (optimiser, learning
    rate, clip parameter, GAE $\lambda$, discount $\gamma$, rollout
    shape, mini-batch size, PPO epochs, hidden-layer width, activation)
    are listed in Appendix~\ref{app:hyperparams}. The DiCo integration
    of Section~\ref{sec:experiments:dico-dispersion} additionally
    specifies the HetControl architecture, $\SND_{\mathrm{des}}{=}0.1$
    set point, $167$-iteration budget, three seeds, $k{=}3$ for the
    $k$-NN graph, Bernoulli inclusion probability $p{=}0.1$, and the
    $128$-env subsample used for the $k$-NN construction. The
    supplementary $n{=}50$ set-point sweep
    (Appendix~\ref{app:dico-n50-sweep}) uses the same $167$-iteration
    budget, three seeds per cell, $p{=}0.1$, and three set points
    $\SND_{\mathrm{des}}\in\{0.12,0.14,0.15\}$.

\item {\bf Experiment statistical significance}
    \item[] Question: Does the paper report error bars suitably and correctly defined or other appropriate information about the statistical significance of the experiments?
    \item[] Answer: \answerYes{}
    \item[] Justification: The unbiasedness diagnostic
    (Section~\ref{sec:experiments:verify}, Appendix~\ref{app:verify})
    reports $95\%$ confidence intervals from $2{,}000$ Bernoulli draws
    per $p$ and the standardised deviation
    $|\mathrm{bias}|/\mathrm{SE}$ as a sharper test. The concentration
    diagnostic reports the empirical violation rate at $\delta{=}0.1$
    over $2{,}000$ draws per cell and the $95$th-percentile deviation.
    The timing experiments
    (Section~\ref{sec:experiments:timing},
    \ref{sec:experiments:scaling}) report mean, median, and standard
    deviation over $20$ timing trials per cell after $3$ warm-ups.
    The DiCo experiment (Section~\ref{sec:experiments:dico-dispersion})
    reports mean $\pm$ standard deviation over three seeds for reward
    and applied SND, and median with IQR for per-call metric cost.
    The $n{=}50$ set-point sweep (Appendix~\ref{app:dico-n50-sweep},
    Table~\ref{tab:dico-n50-sweep}) reports late-window (last
    $50$ iterations) means $\pm$ cross-seed standard error over three
    seeds per set point and the relative tracking error
    $\overline{|\mathrm{applied}\;\SND-\SND_{\mathrm{des}}|}/
    \SND_{\mathrm{des}}$.

\item {\bf Experiments compute resources}
    \item[] Question: For each experiment, does the paper provide sufficient information on the computer resources needed to reproduce the experiments?
    \item[] Answer: \answerYes{}
    \item[] Justification: The $n{\in}\{4,8,16\}$ frozen-checkpoint
    verifications (Section~\ref{sec:experiments:verify},
    Appendix~\ref{app:verify}) run on a single CPU core with the
    reported $20$ trials per cell in minutes. The $n{=}100$ overnight
    PPO run (Section~\ref{sec:experiments:scaling}) fits on one
    RTX~4090, as does the frozen-init $n\in\{100,250,500\}$ GPU timing
    sweep (Section~\ref{sec:experiments:timing}). The DiCo
    closed-loop experiment
    (Section~\ref{sec:experiments:dico-dispersion}) also
    runs on one RTX~4090 at the reported $167$-iteration budget.
    The supplementary $n{=}50$ set-point sweep
    (Appendix~\ref{app:dico-n50-sweep}, nine cells of three seeds
    $\times$ three set points) runs sequentially on one RTX~4090 at
    ${\approx}2$ GPU-hours per cell (${\approx}18$ GPU-hours total).
    Total compute for the experiments reported in this paper is
    under $48$ GPU-hours on one RTX~4090 plus a few CPU-hours for
    the $n{\in}\{4,8,16\}$ benchmarks.

\item {\bf Code of ethics}
    \item[] Question: Does the research conducted in the paper conform, in every respect, with the NeurIPS Code of Ethics?
    \item[] Answer: \answerYes{}
    \item[] Justification: The work is theoretical and empirical
    research on a metric for behavioural diversity in cooperative
    MARL. It does not involve human subjects, sensitive data, or
    release of models with plausible misuse potential. No deviations
    from the NeurIPS Code of Ethics are required.

\item {\bf Broader impacts}
    \item[] Question: Does the paper discuss both potential positive societal impacts and negative societal impacts of the work performed?
    \item[] Answer: \answerNA{}
    \item[] Justification: Graph-SND is a methodological metric
    contribution that reduces the wall-clock cost of measuring
    behavioural diversity in cooperative MARL. It enables neither
    new capabilities nor new deployments; it only makes an existing
    measurement cheaper. We therefore do not identify a direct path
    to positive or negative societal impact.

\item {\bf Safeguards}
    \item[] Question: Does the paper describe safeguards that have been put in place for responsible release of data or models that have a high risk for misuse?
    \item[] Answer: \answerNA{}
    \item[] Justification: No data or models with misuse potential
    are released; the released artefacts are experimental code,
    reproducible CSV/JSON result files, and paper figures.

\item {\bf Licenses for existing assets}
    \item[] Question: Are the creators or original owners of assets used in the paper properly credited and are the license and terms of use explicitly mentioned and properly respected?
    \item[] Answer: \answerYes{}
    \item[] Justification: VMAS \citep{bettini2022vmas} and DiCo
    \citep{bettini2024dico} are credited in the main text and
    released under their original open-source licenses. The vendored
    fork of the DiCo codebase in
    \texttt{ControllingBehavioralDiversity-fork/} retains the
    upstream license, adds a \texttt{GRAPH\_SND\_CHANGES.md} file
    documenting the modifications, and includes the preferred
    citation in \texttt{CITATION.cff}. The original SND paper
    \citep{bettini2025snd} is cited throughout.

\item {\bf New assets}
    \item[] Question: Are new assets introduced in the paper well documented and is the documentation provided alongside the assets?
    \item[] Answer: \answerYes{}
    \item[] Justification: The new assets are the \texttt{graphsnd}
    Python package, the two training scripts, the experiment
    scripts, and the released figure and CSV outputs. The top-level
    \texttt{README.md} documents installation, the four entry-point
    experiments, the overnight training launcher, and the repository
    layout, and cross-references the specific paper sections that
    each asset reproduces.

\item {\bf Crowdsourcing and research with human subjects}
    \item[] Question: For crowdsourcing experiments and research with human subjects, does the paper include the full text of instructions given to participants and screenshots, if applicable, as well as details about compensation?
    \item[] Answer: \answerNA{}
    \item[] Justification: No crowdsourcing or human-subjects
    research is involved.

\item {\bf Institutional review board (IRB) approvals or equivalent for research with human subjects}
    \item[] Question: Does the paper describe potential risks incurred by study participants, whether such risks were disclosed to the subjects, and whether IRB approvals (or equivalent) were obtained?
    \item[] Answer: \answerNA{}
    \item[] Justification: No human-subjects research is involved.

\item {\bf Declaration of LLM usage}
    \item[] Question: Does the paper describe the usage of LLMs if it is an important, original, or non-standard component of the core methods in this research?
    \item[] Answer: \answerNA{}
    \item[] Justification: No LLM is part of the methods, theory,
    or experimental pipeline of this paper. LLMs were used only as
    writing / editing assistance and do not affect the core
    methodology, scientific rigour, or originality of the research.

\end{enumerate}
